\section{The holographic code in detail}\label{app:holocode}

This appendix holds the construction, the methods, the proofs and the remaining results behind Section~\ref{sec:holocode}.

\subsection{Construction and method}\label{sec:holocode-model}

\paragraph{The tile.} A tile is the six-leg perfect tensor of the five-qubit code~\cite{laflamme,bennett1996}:
\[
T_{b\,q_1\dots q_5}=\bra{q_1\dots q_5}E\ket{b}/\sqrt2 .
\]
Here $E$ is the encoding isometry of $[[5,1,3]]$, whose stabilisers are $XZZXI$ and its cyclic shifts, with $\bar X=XXXXX$ and
$\bar Z=ZZZZZ$. The leg $b$ is the logical leg. Every split of the six legs into three and three is unitary~\cite{happy}.

\paragraph{The tree.} The root has its bulk leg $b$ and five legs to children. Every other tile has a bulk leg, a parent leg and four
child legs. The leaves' legs form the boundary, $n=5\cdot4^D$ qubits in ring order, carried by $1+5(4^{D}-1)/3$ tiles. For
$D=1,\dots,5$ this gives $n=20$, 80, 320, 1280 and 5120 boundary qubits on 6, 26, 106, 426 and 1706 tiles. With one reference per
bulk qubit, setting (b) has 6826 qubits at $D=5$.

\paragraph{Orientations.} The children sit on $q_2,\dots,q_5$ in both orientations.
\begin{itemize}[nosep]
\item \emph{Orientation A}, the main construction, is the 25-qubit assembly of Section~\ref{sec:holo}: the parent edge sits on $b$ and
  the bulk on $q_1$.
\item \emph{Orientation H}, the HaPPY convention, puts the bulk on $b$ and the parent edge on $q_1$.
\end{itemize}
Both make every tile a perfect tensor. They differ only in which legs meet across an edge.

\paragraph{Exact computation.} Every tile map is a Clifford isometry, $V\ket v=U\ket{v,0,\dots,0}$. We find $U$ from the dense tensor
and check it against that tensor. The generators are pushed through the tree as signed Pauli strings~\cite{aaronson2004}. In all 20
networks (five depths, two settings, two orientations) they are $N$ independent generators on $N$ qubits. Each one, sign included, is a
$+1$ stabiliser of an independent composition of the same Cliffords by stim~\cite{gidney2021} \chk{holo:H1}. Entropies, in bits,
follow exactly from the generator matrix $G$ over GF(2)~\cite{fattal2004,hamma2005}:
\begin{equation}\label{eq:holocode-rank}
S(A)=\operatorname{rank}\bigl(G|_A\bigr)-|A| .
\end{equation}
At $D=1$ these entropies agree with a dense contraction of the tensors to $2\times10^{-15}$. On a $3\times3$ toric code they match the
dense state vector exactly \chk{holo:H1}, \chk{holo:H6}.

\paragraph{Recovery.} Bulk qubit $t$ is recoverable from region $A$ when $I(R_t\!:\!A)=S(R_t)+S(A)-S(R_tA)=2$ bits. For a stabiliser
state,
\[
I(R_t\!:\!A)=\dim\bigl(V\cap\langle X_{R_t},Z_{R_t}\rangle\bigr),
\]
where $V$ is the reference part of the stabilisers supported on $A\cup R$. One elimination, with columns ordered as complement,
references, $A$, gives $V$ for all bulk qubits at once.

\paragraph{Cuts, wedges and the residual region.}
\begin{itemize}[nosep]
\item The minimal cut $\mathrm{mc}(A)$ assigns every tile to the side of $A$ or of its complement and counts the legs crossed, with unit
  capacity per leg. Dynamic programming on the tree finds the cheapest assignment.
\item In setting (b) a bulk leg is crossed when its tile lies on $A$'s side, because its reference belongs to the complement.
\item The greedy algorithm of Ref.~\cite{happy} absorbs a tile into $A$'s wedge when at least 3 of its 5 non-bulk legs lie in the
  region. We run it to its fixed point.
\item The \emph{residual region} consists of the tiles absorbed neither from $A$ nor from its boundary complement.
\end{itemize}

\subsection{The exceptions to the minimal-cut formula}\label{sec:holocode-rt}

Section~\ref{sec:holocode} gives the exceptions in setting (a) and their cause. Here we characterise them over all setting-(a) cases
\chk{holo:H2}:
\begin{itemize}[nosep]
\item Every exception has a non-empty residual region, so the two greedy cuts never meet.
\item The deficit never exceeded $\lfloor r/2\rfloor$ for $r$ residual tiles. We observed this bound; we did not prove it.
\item Degenerate minimal cuts are not the explanation: 95\,735 exceptions have a unique minimal cut.
\item The argument of Ref.~\cite{happy} for connected regions rests on the greedy cuts meeting. On a tree with fixed bulk they need not
  meet, and then the minimal cut is only an upper bound.
\end{itemize}
One interval is confirmed as an exception in orientation H by dense contraction alone \chk{holo:H1}.

\subsection{Entanglement-wedge reconstruction}\label{sec:holocode-wedge}

\paragraph{The pairs tested.} These are all $2^{20}$ regions, 6 bulk qubits and 2 orientations at $D=1$, plus intervals and random
subsets of 11 sizes at $D=2,\dots,5$ \chk{holo:H3}.

\paragraph{The predictions compared.} In every pair, four predictions agree:
\begin{itemize}[nosep]
\item the exact recoverability;
\item the greedy wedge;
\item the recursion of Proposition~\ref{prop:holocode-erasure};
\item the min-cut prediction $I_{\rm cut}=1+\mathrm{mc}(A)-\mathrm{mc}(A\cup R_t)$.
\end{itemize}
$I_{\rm cut}$ never takes the value 1. Wedge reconstruction~\cite{adh,dong2016} therefore holds in this code.

\paragraph{Two regimes.} Table~\ref{tab:holocode-wedge} shows them.
\begin{itemize}[nosep]
\item An interval holds the bulk geometrically. A deep qubit is in the wedge when its subtree is, and the root needs more than half the
  ring.
\item For a random subset, the root and the upper levels hold nothing until about 77\,\% of the boundary survives, which is $1-p_c$.
  Then they hold nearly everything. Qubits near the leaves are recovered in proportion to their local coverage.
\end{itemize}

\begin{table}[ht]\centering\small
\caption{Fraction of bulk qubits recoverable at $D=5$ (orientation A, setting b), by depth (0 is the root) and by the fraction
$|A|/n$ of the boundary held. Each entry averages 6 interval starts or 6 random subsets. \chk{holo:H3}}\label{tab:holocode-wedge}
\begin{tabular}{lcccccccc}\toprule
& \multicolumn{4}{c}{interval, $|A|/n=$} & \multicolumn{4}{c}{random subset, $|A|/n=$} \\
depth & 0.3 & 0.5 & 0.6 & 0.7 & 0.6 & 0.7 & 0.8 & 0.9 \\\midrule
0 & 0.00 & 0.17 & 0.67 & 1.00 & 0.00 & 0.00 & 1.00 & 1.00 \\
1 & 0.27 & 0.43 & 0.53 & 0.77 & 0.00 & 0.03 & 1.00 & 1.00 \\
3 & 0.29 & 0.50 & 0.60 & 0.70 & 0.07 & 0.44 & 0.99 & 1.00 \\
5 & 0.30 & 0.50 & 0.60 & 0.70 & 0.52 & 0.75 & 0.96 & 1.00 \\\bottomrule
\end{tabular}
\end{table}

\subsection{The erasure recursion}\label{sec:holocode-erasure}

\paragraph{Definitions.} Let $A$ be the boundary qubits that survive erasure. The edge above a tile is \emph{available from below} when
its qubit can be recovered from the part of $A$ below the tile. It is \emph{available from above} when its qubit can be recovered from
the part of $A$ outside the tile's subtree. A surviving boundary leg counts as available.

\begin{proposition}[exact erasure recursion, setting (b); \vd{proven}]\label{prop:holocode-erasure}
\leavevmode
\begin{itemize}[nosep]
\item Seen from either side, every edge is either available or decoupled; that is, its channel to that part of $A$ is constant.
\item An edge is available from below if and only if at least 3 of the 4 legs below its tile are available.
\item An edge is available from above if and only if at least 3 of the parent's other 4 non-bulk legs are available.
\item Bulk qubit $t$ has $I(R_t\!:\!A)=2$ if at least 3 of its tile's 5 non-bulk legs are available, and $I(R_t\!:\!A)=0$ otherwise.
  For the root, this means 3 of its 5 children.
\end{itemize}
\end{proposition}

\noindent\emph{Proof.} Two properties of a perfect tensor are used. Any three legs are a unitary image of the other three, and any
three legs are maximally mixed.
\begin{itemize}[nosep]
\item Suppose at least 3 legs below a tile are available. Decoding them recovers the remaining legs, of which there are at most three:
  the edge, the bulk leg and at most one child.
\item Suppose at most 2 legs are available. Then the edge and those legs form three maximally mixed legs, so $I=1+2-3=0$.
\item By induction on the height, an unavailable subtree is a constant channel. After local decoding, an available subtree is
  equivalent to its edge qubit.
\end{itemize}
The same argument runs from above, and with the bulk leg as the target. \hfill$\square$

\paragraph{The probabilities.} The availabilities in Proposition~\ref{prop:holocode-erasure} depend on disjoint sets of boundary
qubits, so they are independent. This gives Eq.~\eqref{eq:holocode-rec}, with the binomial step
$a_h=P(\mathrm{Bin}(4,a_{h-1})\ge3)$. Analogous closed forms hold at every depth. At $a^*$ the map has slope
$f'(a^*)=12a^{*2}(1-a^*)=1.643>1$. Its fixed point is therefore unstable, and the transition window narrows as $1.643^{-D}$
\chk{holo:H4}.

\begin{table}[ht]\centering\small
\caption{Probability that the root is recoverable after erasure at rate $p$: the Monte Carlo estimate, with the recursion
\eqref{eq:holocode-rec} in brackets. There are 4000, 2000, 1500, 1000 and 400 samples per $p$ for $D=1,\dots,5$. The row $0.2324$
is $p_c$. \chk{holo:H4}}\label{tab:holocode-mc}
\begin{tabular}{lccccc}\toprule
$p$ & $D=1$ & $D=2$ & $D=3$ & $D=4$ & $D=5$ \\\midrule
0.22 & 0.938 (0.932) & 0.946 (0.942) & 0.958 (0.956) & 0.976 (0.973) & 0.990 (0.988) \\
0.2324 & 0.913 (0.914) & 0.909 (0.914) & 0.907 (0.914) & 0.905 (0.914) & 0.912 (0.914) \\
0.25 & 0.886 (0.884) & 0.854 (0.860) & 0.821 (0.815) & 0.704 (0.722) & 0.517 (0.529) \\
0.28 & 0.818 (0.819) & 0.713 (0.730) & 0.549 (0.546) & 0.247 (0.236) & 0.020 (0.017) \\\bottomrule
\end{tabular}
\end{table}

\paragraph{Checks of the recursion.}
\begin{itemize}[nosep]
\item At $D=1$ the recursion equals, to $3\times10^{-13}$, the exact recovery polynomial built from all $2^{20}$ erasure patterns.
\item The Monte Carlo samples of Section~\ref{sec:holocode} cover orientation A at $D=1,\dots,5$, at 11 values of $p$ from 0.05 to
  0.5. A further check in orientation H at $D=5$ agrees in the same way.
\item The Monte Carlo fractions match the probabilities at every point of $D$, $p$ and depth, with largest $|z|=2.55$.
\item The Monte Carlo curves for $D=4$ and $D=5$ cross at $p=0.2331$, within the noise of 400 samples \chk{holo:H4}.
\end{itemize}

\subsection{The witnesses in detail}\label{sec:holocode-witness}

\paragraph{Contents.} The six bulk contents of Section~\ref{sec:holocode} are five random products of cardinal states and one with 400
bulk Bell pairs.

\paragraph{The two witnesses.} Witness I reads the basis $g_1,g_1g_2,g_2g_3,\dots$ of the stabiliser group, not the generators
themselves, and decodes the result. In the noisy runs, both witnesses also match the commutation pattern of the error with the
generators, which is computed separately.

\paragraph{Reading the content.} After the content of the 86 read qubits is read, the other 1620 bulk qubits keep fidelity 1
\chk{holo:H5}.

\paragraph{One tile.} We enumerated every Pauli error on the five qubits of a tile. The syndrome never depends on the content, and lookup
correction restores every single-qubit error. By weight $0,\dots,5$, the correctable errors number $1,15,0,60,135,45$. Under
depolarising noise of rate $p$, the logical failure probability is therefore exactly
\begin{equation}\label{eq:holocode-pfail}
P_{\rm fail}=1-\Bigl[(1-p)^5+5p(1-p)^4+\tfrac{20}{9}p^3(1-p)^2+\tfrac53p^4(1-p)+\tfrac{5}{27}p^5\Bigr]=10p^2+O(p^3).
\end{equation}
\begin{itemize}[nosep]
\item This lies below the closed form $1-(1-p)^5-5p(1-p)^4$ for a distance-3 code that corrects only weight-1 errors. The two give
  $0.0795$ against $0.0815$ at $p=0.1$, and $0.249$ against $0.263$ at $p=0.2$.
\item The difference comes from errors of weight 3 to 5 that equal a correctable error times a stabiliser.
\item The pseudo-threshold, where $P_{\rm fail}=p$, is $p=0.1376$.
\item For Haar-random contents the syndrome distribution is independent of the content to machine precision. Reading the content
  leaves fidelity exactly $\tfrac12$ \chk{holo:H5}.
\end{itemize}

\subsection{The classical centre}\label{sec:holocode-centre}

Section~\ref{sec:holo} closes on Harlow's formula~\cite{harlow},
$S(\rho_R)=\Tr(\rho\,\mathcal L_R)+S_{\rm bulk}(\rho,\mathcal M_R)$. In it the area operator $\mathcal L_R$ lies in the centre of the
reconstructable algebra $\mathcal M_R$. In setting (b) we can read the formula off directly, in three steps.
\begin{itemize}[nosep]
\item The cut function is submodular, so minimal cuts are closed under union. The largest $A$-side among them is therefore itself a
  minimal cut.
\item $I_{\rm cut}(t)=2$ holds exactly when tile $t$ lies on $A$'s side of some minimal cut. So this largest side is the set of tiles
  with $I_{\rm cut}=2$. Because $I_{\rm cut}=I$ in every pair, it is the wedge $W_A$.
\item The cut crosses one bulk leg per tile of $W_A$ and $|\gamma_A|$ non-bulk legs. Hence
\end{itemize}
\begin{equation}\label{eq:holocode-qrt}
S(A)=|\gamma_A|+|W_A|\qquad\text{(setting b)}.
\end{equation}

\paragraph{An empty residual region.} Every bulk qubit then lies either in $W_A$ or in the wedge of the complement. Operators on $A$
commute with operators on the complement, so $\mathcal M_A$ is exactly the algebra of the qubits in $W_A$. Equation~\eqref{eq:holocode-qrt}
is then Harlow's formula, with $S_{\rm bulk}=|W_A|$ and $\Tr(\rho\,\mathcal L_A)=|\gamma_A|$, which is a number. On a fixed tree the area
can only be a number. The assembly realises one fixed-area sector, as Lemma~\ref{lem:classical} requires, and it cannot represent a
superposition of areas. That \vd{holds} by construction. The superposition of geometries stays \vd{untested}.

\paragraph{Setting (a).} With an empty residual region and a pure bulk, Harlow's formula gives $S=|\gamma_A|$, and we find
$S=\mathrm{mc}$ in every such case.

\paragraph{A non-empty residual region.} A residual tile is a bulk qubit that neither side can reconstruct alone. Complementary
recovery of whole bulk qubits then fails, and Harlow's theorem in this form does not apply. The deficits have the tensor-network cause
given in Section~\ref{sec:holocode}. Whether a non-trivial centre of the reconstructable algebra also describes them is
\vd{untested}: $I(R_t\!:\!A)$ takes only the values 0 and 2, and the script constructs no centre.

\subsection{A control: the toric code}\label{sec:holocode-toric}

We built the ground state of the toric code~\cite{kitaev2003} on $L\times L$ tori for $L=8$, 16, 24, 32, 40 and 45, which reaches 4050
qubits. The generators are the stars, the plaquettes and two non-contractible $Z$ loops; their GF(2) rank is $2L^2$. The regions are
sets of edges in disks or annuli about a generic centre \chk{holo:H6}. We applied two constructions.
\begin{itemize}[nosep]
\item \emph{Kitaev and Preskill}~\cite{kitaev2006}, $-\gamma=S_A+S_B+S_C-S_{AB}-S_{BC}-S_{CA}+S_{ABC}$, on 33 disks of radius 2 to 18.
  Each disk is cut into three sectors.
\item \emph{Levin and Wen}~\cite{levin2006}, $-2\gamma=(S_1-S_2)-(S_3-S_4)$, on 20 annuli, from radii 2--4.5 to 10--19.5. Region 1 is
  the annulus, regions 2 and 3 have one gap each, and region 4 has both gaps.
\end{itemize}
Every case gives $\gamma=1$ bit $=\ln2$, the value for $\mathbb Z_2$ topological order, whatever the region size and $L$ \tagP. The same
machinery that finds the deficits of Table~\ref{tab:holocode-rt} reproduces this universal constant exactly.

\subsection{Choices and caveats}\label{sec:holocode-caveats}

\paragraph{Choices made after seeing the data.}
\begin{itemize}[nosep]
\item Orientation H was added after the exhaustive run at $D=1$ showed interval exceptions in the HaPPY convention.
\item A first characterisation, that exceptions need a degenerate minimal cut, failed on 95\,735 regions. The residual region and the
  observed $\lfloor r/2\rfloor$ bound replaced it.
\item The first witness run at scale failed: shortening the pushed generators had mixed reference factors into the code stabilisers.
  Shortening now uses only reference-free generators, and no entropy changes.
\item The Monte Carlo samples at $D=5$ were raised from 160 to 400.
\item Two floating-point tolerances were widened tenfold, to $10^{-11}$ and $10^{-13}$, after rounding exceeded them.
\item A counting error in the script had added the setting-(b) sets at $D=1$ to the setting-(a) total. It is fixed. The checks and every
  other number are unchanged.
\item The region sets were sized for runtime.
\end{itemize}

\paragraph{Caveats.}
\begin{itemize}[nosep]
\item This is kinematics on a fixed background.
\item The states are stabiliser states. Non-stabiliser contents were tested on one tile only.
\item The noise is Pauli or erasure noise, and syndrome extraction is noiseless.
\item The $\lfloor r/2\rfloor$ bound is observed in every setting-(a) case, not proved, and the regions at $D\ge2$ are samples.
\item The recursion is exact only in setting (b).
\item A tree is not our space (Section~\ref{sec:holo}).
\end{itemize}
